%% Clear sans, courier

% Klasa:
\documentclass[12pt]{mwart}

% Kodowanie:
\usepackage[utf8]{inputenc}

% Czcionki:
\usepackage[sfdefault]{ClearSans} %% option 'sfdefault' activates Clear Sans as the default text font
\usepackage[T1]{fontenc}
\usepackage{courier}
\usepackage[T1]{polski}

% Geometria i style
%\usepackage{showframe} % debug
\usepackage{geometry}
\geometry{a4paper,twoside=true,top=20mm,bottom=30mm,inner=35mm,outer=20mm}
\pagestyle{plain}
\usepackage{float}
\usepackage{afterpage}

% Grafika
\usepackage{graphicx}

% kody
\usepackage{listings}

% matma
\usepackage{amssymb,amsmath}

% Numerowanie obiektow wedlug rozdzialow
\usepackage{chngcntr}
\counterwithin{figure}{section}

% Linki
\usepackage{url}

% Dane:
\title{Uniwersalny inteligentny sterownik silników do zastosowań przemysłowych}
\author{Łukasz Przeniosło}

% Zmienne srodowiskowe:
\newcommand{\paperTitlePl}{UNIWERSALNY INTELIGENTNY STEROWNIK SILNIKÓW DO ZASTOSOWAŃ PRZEMYSŁOWYCH}
\newcommand{\paperTitleEn}{UNIVERSAL INTELIGENT MOTOR DRIVER FOR INDUSTRIAL APPLICATIONS}
\newcommand{\topicChooseDate}{01.08.2015 r.}
\newcommand{\topicAcceptationDate}{01.06.2016 r.}
\newcommand{\tocLessLeft}[1]{\textsc{\textbf{\large \flushleft{#1}}} \vspace{3mm} \\ \indent }
\newcommand{\tocLessCenter}[1]{\textsc{\textbf{\large \centerline{#1}}} \vspace{3mm} \\ \indent }
\newcommand{\insertImg}[3]{
\begin{figure}[h]
	\centering
	\includegraphics[height=50mm]{#1}
	\caption{#2} \flushleft{Źródło: Na podstawie \cite{#3}}
\end{figure}
}

% Indeksowanie w pdfie wyjsciowym
\usepackage{hyperref}

% START:
\begin{document}
\thispagestyle{empty} % usun nr strony ze strony tytulowej

\begin{figure}[h]
	\centering
	\includegraphics[height=18mm]{grafiki/logo_zut.pdf}
	\hspace{5mm}
	\includegraphics[height=18mm]{grafiki/logo_zut_we.pdf} \\
	\vspace{20mm}


	{\bf \large Łukasz Przeniosło} \\ 
	\vspace{5mm}

	nr albumu: 21789 \\
	kierunek studiów: Elektrotechnika \\
	specjalność: urządzenia i instalacje elektryczne \\
	forma studiów: studia niestacjonarne \\
	\vspace{20mm}

	\textsc{\textbf{\large \paperTitlePl}} \\
	\vspace{5mm}

	{\bf \large \paperTitleEn} \\
	\vspace{45mm}
	
	
	Praca dyplomowa magisterska \\
	napisana pod kierunkiem: \\
	\vspace{2mm}
	
	{\bf \large dr inż. Michała Bonisławskiego} \\
	\vspace{2mm}
	
	Katedra Elektroenergetyki i Napędów Elektrycznych
	\vspace{5mm}
	
\end{figure}

\begin{flushleft}
	{\footnotesize Data wydania tematu pracy: \topicChooseDate} \\
	{\footnotesize Data złożenia tematu pracy: \topicAcceptationDate} 
\end{flushleft}

\vspace{\fill \centerline{Szczecin 2016}}
\clearpage

\afterpage{\null \thispagestyle{empty} \newpage}
\clearpage

% Oswiadczenie autora
\thispagestyle{empty}

\tocLessCenter{OŚWIADCZENIE AUTORA PRACY DYPLOMOWEJ} 

Oświadczam, że praca dyplomowa inżynierska/magisterska pt. \\
 „\paperTitlePl” \\
napisana pod kierunkiem: dr. inż. Michała Bonisławskiego \\
jest w całości moim samodzielnym autorskim opracowaniem sporządzonym przy wykorzystaniu wykazanej w pracy literatury przedmiotu i materiałów źródłowych.
Złożona w dziekanacie Wydziału Elektrycznego treść mojej pracy dyplomowej w formie elektronicznej jest zgodna z treścią w formie pisemnej. \\

Oświadczam ponadto, że złożona w dziekanacie praca dyplomowa ani jej fragmenty nie
były wcześniej przedmiotem procedur procesu dyplomowania związanych z uzyskaniem
tytułu zawodowego w uczelniach wyższych.
\vspace{30mm}


\begin{flushright}
	\begin{tabular}{c}
  		............................................\\
  		podpis dyplomanta
	\end{tabular}
\end{flushright}

\vspace{5mm}

\begin{flushleft}
	Szczecin, dn. \topicAcceptationDate{}
\end{flushleft}


\clearpage


% Streszczenie
\thispagestyle{empty}

\tocLessLeft{Streszczenie pracy}
Praca przedstawia projekt systemu sterowniczego do różnorodnych silników elektrycznych, na który składają się między innymi: Zbudowane w oparciu o 32 bitowy mikrokontroler urządzenie do którego podłączany jest silnik i peryferia (sterownik) oraz aplikacja PC-owa służąca do nastaw niezbędnych do poprawnego działania układu parametrów brzegowych oraz badania pracy zasilanego przez sterownik silnika. Rozwiązanie ma zostać zastosowane do sterowania elementami ruchomymi w maszynach typu {\it Pick and Place}, drukarkach szablonowych, drukarkach 3D oraz we wszelkiego rodzaju konwojerach. W części opisowej dokumentu, została zaprezentowana niezbędna teoria nt. silników elektrycznych oraz sterowania. Ponadto, zostały porównane komercyjne rozwiązania sterowników różnej klasy, które są dostępne na rynku.

\tocLessLeft{Słowa kluczowe}
Sterowanie, silnik, elektronika mocy, oprogramowanie

\tocLessLeft{Abstract}
The paper presents a project of a control system, that can be used with a various types of electric motors. Among others the main components of the system are: built upon a 32 bit microcontroller device, to which the motors and peripherals are connected (the driver) and a PC application used for defining critical operating parameters and for gathering connected motor work diagnostics . This application is to be used with Pick and Place machines, stencil printers, 3D printers and all sort of conveyors. In the generic part of the text, the essential theory of electric motors and control has been described. Also, commercially available on the market  motor drivers of various class has been compared.

\tocLessLeft{Key words}
Controll, motor, power electronics, software


\clearpage


% Spis tresci
\tableofcontents
\clearpage



% Wstep
\section*{Wprowadzenie}

Chcąc zaprojektować system napędowy dla określonej aplikacji należy wziąć pod uwagę bardzo wiele czynników. Najważniejsze z nich to: koszt i dostępność rozwiązań gotowych na rynku. Ceny rozwiązań przemysłowych bardzo rzadko mają odzwierciedlenie w zastosowanym w ich strukturze {\it hardwarze}, jest to raczej cena za {\it know-how} producenta, wiedzę ekspercką, wsparcie techniczne oraz dziesiątki roboczogodzin programistów i elektroników tworzących produkt. W przypadku posiadania dużych środków takie rozwiązanie wydaje się być najbardziej odpowiednie. Ma ono jednak jedną znaczącą wadę (pomijając cenę oczywiście). Pomimo tego że producenci sterowników do silników prześcigają się w ilości posiadanych cech, czy technologii, to korzystając z takiego urządzenia nijako jesteśmy zmuszeni z korzystania z zamkniętego, ograniczonego systemu. System taki może lepiej lub gorzej wpasowywać się w ramy naszej aplikacji. Rzadko kiedy wpasowuje się idealnie, a najczęściej brakuje mu jakichś właściwości, lub jest przewymiarowany. \

Drugą opcją jest zakup taniego zintegrowanego chipu, który po raz kolejny w zależności od wyłożonych środków, (w o wiele mniejszej skali) zawiera w sobie mniej lub więcej logiki i czasami końcówki mocy. Taki układ można także zbudować z elementów dyskretnych, co nie obniża znacząco kosztów, lecz daje możliwość większej kontroli nad parametrami układu. Wady takiego rozwiązania są takie, że zintegrowane sterowniki są projektowane pod zastosowanie z konkretnym rodzajem silnika (np tylko do silników krokowych) i służą do sterowania silnikami małej mocy jeśli mają zintegrowane końcówki. W przypadku stosowania układów tego typu oraz dyskretnych, projektant jest zmuszony zaprojektować obwód drukowany oraz napisać bardzo dużą ilość kodu zawierającego wiele skomplikowanych procedur, które muszą działać w czasie rzeczywistym. \

W przypadku, kiedy ilość produkowanych maszyn współpracującymi ze sterownikami do silników jest duża, posiadając odpowiednią wiedzę i umiejętności, bardziej opłaca się zaprojektować rozwiązanie własne, "od zera". Ponadto, proces tworzenia tego typu systemu dostarcza ogromną ilość wiedzy praktycznej.
  
\subsection*{Cel pracy}

Zbudować sterownik do sterowania różnego typu silników (krokowych, {\it BLDC}, liniowych i {\it DC}), z możliwością dostosowywania parametrów, używając aplikacji PC. Aplikacja ma pozwalać na strojenie działających w sterowniku regulatorów, wybierać tryby pracy i badać działanie systemu w czasie rzeczywistym. Urządzenie ma zostać użyte w maszynach przemysłowych służących do produkcji urządzeń elektronicznych.

\subsection*{Zakres pracy}

Przegląd i analiza rozwiązań technologicznych, układów elektronicznych i oprogramowania stosowanych w sterownikach silników oferowanych przez wybranych producentów. Projekt uniwersalnego, kompaktowego sterownika do silników {\it DC} (prądu stałego), {\it BLDC (Brushless DC Motor)}, krokowych oraz liniowych {\it VCM (Voice Coil Motor}, inaczej {\it VCA - Voice Coil Actuator}) wraz z opisem sposobów regulacji i algorytmów zawierających elementy {\it AI}, które zostaną zaimplementowane w nowoczesnym 32 bitowym mikrokontrolerze z rodziny {\it ARM}. Opracowanie aplikacji na PC umożliwiającej dostrajanie parametrów dla konkretnego silnika oraz analizowanie jego pracy. Wykonanie, uruchomienie i badania laboratoryjne prototypu urządzenia. Prezentacja przykładowego zastosowania sterownika działającego w rzeczywistym układzie.

\clearpage


%Start czesci teoretycznej
\section{Elektromagnetyzm}

Aby zrozumieć zasadę działania większości silników elektrycznych, należy najpierw zaznajomić się z podstawowymi jednostkami występującymi w elektromagnetyzmie, dzięki którym silniki działają. Rozdział został opracowany na podstawie źródła \cite{jaszczuk}

\subsection{Natężenie pola magnetycznego}

\insertImg{grafiki/right_hand_rule.pdf}
	   {Reguła prawej dłoni - sposób wyznaczania zwrotu linii sił pola magnetycznego wokół przewodnika w którym jest prąd elektryczny}
	   {wikiRegPrawDlon}








% Opis silnikow (podrozdzialy beda includowane w rozdzialach)
\section{Podstawowe informacje o silnikach}

Rozdział ten przedstawia podział maszyn elektrycznych ze względu na ich różne parametry, oraz opis tych silników, które są (i mogą być) wspierane przez skonstruowany sterownik. 
Rozdział został opracowany głównie na podstawie źródeł \cite{przepiorkowski} i \cite{jones}.

\subsection{Podział silników elektrycznych}

Istnieje bardzo wiele rodzajów silników, a ponadto cały czas powstają nowe konstrukcje, które trudno jest sklasyfikować przy pomocy klasycznych definicji. Podział na silniki AC i DC nie wystarczy-- w przypadku zastosowania elektronicznego modułu sterującego rodzaj zasilania silnika nie ma tak dużego znaczenia. Nie można także jednoznacznie stwierdzić, czy maszyna z wirującymi magnesami to silnik prądu stałego, czy przemiennego. Biorąc to pod uwagę, silniki elektryczne można podzielić na trzy podstawowe grupy.

\tocLessLeftNorm{Silniki prądu stałego -- DC} 

Zaliczane są tu silniki komutatorowe:
\begin{easylist}
	& z magnesem trwałym (PMDC),
	& szeregowe,
	& bocznikowe,
	& szeregowo--bocznikowe.
\end{easylist} 

\tocLessLeftNorm{Silniki prądu przemiennego -- AC}

Do tej grupy  zaliczane są silniki zasilane bezpośrednio z sieci prądu przemiennego. W tym przypadku podział jest bardziej skomplikowany, ze względu na inne kryteria:
\begin{easylist}
	& \underline{Silniki asynchroniczne} -- charakteryzuje je tzw. ,,poślizg''. Rotor takiego silnika wiruje z niższą prędkością od prędkości wirowania pola magnetycznego. Najpopularniejszymi silnikami z tej grupy są silniki klatkowe.
	& \underline{Silniki synchroniczne} -- rotor wiruje z prędkością równą prędkości wirowania pola magnetycznego. Mogą być jednofazowe lub trójfazowe. Ze względu na konstrukcję, silniki te dzielą się na reluktancyjne lub z wirującym magnesem.
	& \underline{Silniki komutatorowe uniwersalne} -- komutatorowe silniki szeregowe przystosowane konstrukcyjnie do zasilania z jednej fazy prądu przemiennego.
\end{easylist} 

\tocLessLeftNorm{Silniki z komutacją elektroniczną}

Do tej grupy zalicza się silniki, których działanie nie było by możliwe bez elektronicznego układu sterującego. Najpopularniejsze z nich to:
\begin{easylist}
	& \underline{Silniki krokowe} -- ze względu na konstrukcję wyróżnia się silniki:
		&& o zmiennej reluktancji,
		&& z magnesem trwałym,
		&& hybrydowe.
\end{easylist} 
Ze względu na sposób sterowania uzwojeniami:
\begin{easylist}
		&& silniki unipolarne,
		&& silniki bipolarne.
\end{easylist} 
	Daje to sześć różnych kombinacji konstrukcyjnych silników krokowych.
	
\insertImg{grafiki/stepper_nema17.png}
	{Przykładowy silnik krokowy hybrydowy ST4118 w standardzie Nema17, produkowany przez firmę Nanotec}
	{nanotec}	
	
Nema ({\em National Electrical Manufacturers Association}) to stowarzyszenie producentów urządzeń elektrycznych, którego siedziba mieści się w USA. Standard Nema określa między innymi wymiary silników krokowych. Numer 17 oznacza szerokość i wysokość silnika 1.7 x 1.7 cala (43.2 x 43.2 mm).

\insertImg{grafiki/bldc_db42.png}
	{Przykładowy silnik BLDC (model DB42), produkowany przez firmę Nanotec}
	{nanotec}

\begin{easylist}
	& \underline{Silniki bezszczotkowe z wirującym magnesem (BLDC)} -- w literaturze angielskiej silniki te są opisywane jako prądu stałego ({\em BrushLess DC}). W przypadku silników dwubiegunowych z czujnikami Halla (np. tego typu, które są używane w małych wentylatorach do chłodzenia radiatorów) taki opis może być uzasadniony. Nie można jednak jednoznacznie zaliczyć do grupy DC silników wielobiegunowych trójfazowych z wirującym magnesem \cite{przepiorkowski}. Mylne może być także to, że w przypadku silników z wbudowanym sterownikiem nie różnią się one z punktu użytkownika od silnika komutatorowego z magnesem trwałym. Maszyny BLDC charakteryzują się bardzo wysoką prędkością obrotową, w szczególności w porównaniu do silników krokowych.	
	
	& \underline{Silniki liniowe} -- wynikiem działania takiej maszyny jest przemieszczenie liniowe, a nie ruch obrotowy. Konstrukcja takiego silnika w dużym uproszczeniu polega na ,,rozwinięciu'' wirnika i stojana na płaszczyźnie. Nazwa silnik liniowy lub {\em Linear Actuator} bardzo często dotyczy zwykłego silnika krokowego z wbudowaną przekładnią śrubową zamieniającą ruch obrotowy na posuwisto--zwrotny, co bardzo często wprowadza niejasności i wątpliwości podczas projektowania systemów napędowych.
	
	\insertImg{grafiki/linear_actuator_l35.png}
		{Przykładowy silnik krokowy hybrydowy L35 z przekładnią śrubową, produkowany przez firmę Nanotec}
		{nanotec}
	
	& \underline{Silniki liniowe VCM} -- zasada działania silnika VCM ({\em Voice Coil Motor}) lub VCA ({\em Voice Coil Actuator}) jest bardzo zbliżona do zasady działania cewki głośnikowej w szczelinie magnesu. Przesunięcie liniowe tych silników nie przekracza najczęściej 50 mm.		
	
\end{easylist} 

\insertImg{grafiki/vcm_moticont.png}
	{Przykładowy silnik liniowy VCM (model GVCM-051-051-01), produkowany przez firmę Moticont}
	{moticont}
	
\subsection{Parametry silników}

Podstawowe parametry silnika są czasami podawane na jego tablicy znamionowej. Pozostałe należy odszukać w jego karcie katalogowej lub oszacować według właściwości poszczególnych rodzajów silników.
\\

\begin{easylist}

	& Moc znamionowa: podawana jest moc elektryczna w watach pobierana przez silnik przy pracy z normalną prędkością obrotową. Moc mechaniczna jest mniejsza i zależy od sprawności silnika (typowo od 40 do 80 \%).
	\\
	& Napięcie zasilania: znamionowa wartość napięcia zasilającego (stałego lub zmiennego), przy której określane są pozostałe parametry. Najczęściej silniki mogą być zasilane napięciem niższym. Nie należy stosować napięć dużo wyższych od znamionowego (maks. +15 \%).
	\\
	& Moment obrotowy: parametr ten jest podawany w niutonometrach $ (Nm) $. Typowe wartości dla małych silników: od 0.01 do 50 Nm. Dla większości rodzajów silników występuje zależność momentu obrotowego i natężenia prądu pobieranego przez maszynę.
	\\
	& Moment rozruchowy: ważny parametr informujący o tym czy silnik jest w stanie wystartować pod obciążeniem. Podawany w Nm lub w procentach momentu obrotowego. W zależności od rodzaju silnika moment może być mały (do 150 \%), średni (150 - 250 \%) lub duży (> 250 \%).
	\\
	& Obroty znamionowe: wyrażane w obrotach na minutę ($ RPM $ -- {\em Revolutions per minute}. Zawsze podawane są przy znamionowym obciążeniu i napięciu zasilania lub podana jest charakterystyka obrotów w funkcji obciążenia. Obroty silników zawierają się w granicach od 100 do 100000 RPM, ale najczęściej spotykane wartości to od 1000 do 4000 RPM. Dla maszyn prądu przemiennego podawane są obroty przy określonej częstotliwości napięcia sieciowego (50 lub 60 Hz). Dla silników krokowych obroty znamionowe nie są podawane. Zamiast tego występują obroty maksymalne lub dopuszczalna częstotliwość impulsów.
	\\
	& Prąd znamionowy i rozruchowy: pobór prądu w normalnych warunkach pracy i w czasie rozpędzania silnika. Ten drugi może być nawet 2 - 8 krotnie większy od znamionowego (największe prądy rozruchowe mają silniki indukcyjne). Prąd rozruchowy nie występuje w silnikach sterowanych impulsowo.
	\\
	& Kierunek obrotów: parametr dotyczący przede wszystkim silników komutatorowych szeregowych i niektórych silników indukcyjnych jednofazowych. Dla tych silników kierunek obrotów jest wymuszony przez konstrukcję silnika. \newline 
	W pozostałych maszynach kierunek zależy od polaryzacji przyłożonego do ich zacisków napięcia lub w przypadku maszyn ze sterowaniem impulsowym-- od kolejności podawanie impulsów na poszczególne uzwojenia.
	\\
	& Konstrukcja mechaniczna: wymiary, ciężar, średnica wału, mocowania, rodzaj łożysk wirnika itp.
	\\
	& Konstrukcja obudowy, chłodzenie: spotykane są silniki o konstrukcji otwartej, w której uzwojenia są dostępne z zewnątrz lub zamkniętej w której silnik jest całkowicie obudowany. W silnikach otwartych często wymagane jest chłodzenie strumieniem powietrza-- np. silniki odkurzaczy, czy wiertarek.  

\end{easylist} 

\subsection{Silnik liniowy (VCM)}

Pierwszym z silników obsługiwanych przez projektowany sterownik jest silnik liniowy VCM. Jego konstrukcja jest bardzo prosta, aczkolwiek sam proces technologiczny wymagany do jego wytworzenia już nie. Z tego właśnie powodu VCM nawet bardzo małej mocy są drogie (Cena od \$100 za pojedyncze sztuki). Ciężko jest tak naprawdę stwierdzić czy jest to silnik czy siłownik, bo zakres jego ruchu nie przekracza zwykle kilku centymetrów.

\insertImg{grafiki/vcm_budowa.png}
	{Budowa cylindrycznego silnika VCM}
	{przepiorkowski}

Nazwa silnika pochodzi od jego zasady działania- identycznie jak w zwykłym głośniku jest to cewka poruszająca się w szczelinie magnesu. Zgodnie z regułą Lorentza:

\begin{defn}
	Jeżeli przez cewkę znajdującą się w polu magnetycznym przepływa prąd, to na cewkę działa siła proporcjonalna do natężenia prądu, a napięcie samoindukcji jest proporcjonalne do szybkości poruszania się cewki. \cite{przepiorkowski}.
\end{defn}

Podstawowe części składowe silnika VCM to:
\begin{easylist}
	& dwa magnesy trwałe ułożone tak, aby były skierowane tym samym biegunem w stronę cewki,
	& cewka poruszająca się w szczelinie między magnesami,
	& obudowa z miękkiego żelaza zamykająca obwód magnetyczny.
	\\
\end{easylist}

Zmieniając natężenie i kierunek przepływającego prądu, możliwe jest bardzo precyzyjne sterowanie położeniem cewki. Z tego powodu układ elektroniczny sterujący silnikiem nie wymaga wielu komponentów, a algorytm sterowania nie jest skomplikowany to VCM ma jednak jeszcze jedną zaletę, która sprawia że jest on wybierany do niektórych aplikacji zamiast np. tańszego silnika krokowego z przekładnią śrubową- pozwala wykonywać bardzo dynamiczne ruchy i posiada wysoką zwrotność. ,,Liniowy'' silnik krokowy o podobnej mocy nawet z zastosowaniem śruby o dużym skoku, nie będzie w stanie osiągnąć podobnych prędkości przy zachowaniu korzystnego momentu. Brak konieczności zamiany ruchu obrotowego na posuwisto-zwrotny eliminuje błędy pozycji (które mogą wynikać np. z luzów w przekładni). Silniki VCM stosowane są wszędzie tam, gdzie potrzebne jest dokładne i szybkie ustawienie pozycji, np. w precyzyjnych urządzeniach mechanicznych i optycznych. \\

Aby efektywnie i bezpiecznie korzystać z silników VCM w danej aplikacji, należy mieć na uwadze kilka ważnych cech tych urządzeń:

\begin{easylist}
	& Korpus cewki nie jest w żaden sposób przymocowany do obudowy silnika. Oznacza to że to projektant musi zadbać o to, aby korpus wsuwał i wysuwał się z cylindra równolegle, stosując odpowiednie mocowanie mechaniczne obu części. Nie zapewniwszy tego, uzwojenie silnika jest narażone na starcie się lakieru izolacyjnego, poprzez ocieranie się o obudowę. 
	
	& W zależności od aktualnego położenia (wysunięcia) korpusu, zależność między siłą a przepływającym przez uzwojenie prądem jest zmienna.
	
	\insertImgSetSize{grafiki/vcm_force_current_plot.eps}
		{80}
		{Przykładowy przebieg siły $ [\frac{N}{A}] $ pchającej/ ciągnącej w funkcji wysunięcia korpusu dla silnika liniowego VCM 019-048-02, firmy Moticont.}
		{moticont}
	
Graf \ref{grafiki/vcm_force_current_plot.eps} prezentuje zależność między siłą silnika a położeniem jego cewki. Jak widać największa i stała siła ($ 2 \frac{N}{A} $) występuje dopiero kiedy korpus jest trochę wysunięty i zaczyna lawinowo spadać przy końcowym wysuwie. Jeżeli układ mechaniczny, którego częścią jest silnik, nie zostanie wyposażony w odpowiednią blokadę, korpus może wypaść z cylindra ponieważ silnik nie będzie już w stanie go utrzymać.

	& Producent w karcie katalogowej silnika podaje parametr maksymalnej ciągłej mocy ({\em Max continuous power}) i odpowiada ona zazwyczaj prądowi, który płynie przez cewkę w momencie niepełnego wysunięcia (należy oczywiście wziąć pod uwagę fizyczne obciążenie części ruchomej- masę obudowy lub korpusu, w zależności od tego która z części jest przymocowana, a która jest w ruchu). Przy dużym wysunięciu zależność siły do płynącego przez cewkę prądu drastycznie spada i utrzymując korpus długo w takiej pozycji można w bardzo łatwy sposób spalić uzwojenie silnika. \\
	
\end{easylist}

Przy zastosowaniu odpowiedniej konstrukcji układu napędowego (zachowana równoległość części mechanicznych) i sterowania zapewniającego ochronę przed zbyt dużym prądem płynącym przez cewkę, silnik VCM z czasem nie ulega właściwie żadnej degradacji i długość jego działania jest ograniczona żywotnością części mechanicznych układu.










\subsection{Silnik krokowy}

Silniki krokowe są maszynami elektrycznymi bez komutatorowymi. Zazwyczaj wszystkie uzwojenia silnika są częścią stojana, podczas gdy rotor jest magnesem trwałym lub w przypadku silników VRM ({\em Variable Reluctance Motor}) zębatym blokiem z materiału magnetycznie miękkiego (patrz sekcja \ref{subsec:materialy_fer}). Komutacja (czyli zmiana kierunku/ uzwojenia w którym następuje przepływ prądu elektrycznego) musi być wymuszona poprzez zewnętrzne urządzenie sterujące maszyną. Takie sterowniki są zazwyczaj projektowane w taki sposób, aby zapewnić możliwość utrzymania wirnika w niemal dowolnej pozycji kątowej, oraz obracać nim w obu kierunkach. 

\subsubsection{Rodzaje silników krokowych}

Silniki krokowe ze względu na konstrukcję można podzielić na trzy główne grupy: 
\begin{easylist}
	& silniki o zmiennej reluktancji,
	& z magnesem trwałym,
	& hybrydowe.
\end{easylist}

\insertImg{grafiki/reluctance_moment.png}
		{Powstawanie reluktancyjnego momentu obrotowego.}
		{przepiorkowski}

Reluktancja jest parametrem analogicznym do rezystancji elektrycznej, lecz odniesionym do strumienia magnetycznego. Jest to inaczej ,,rezystancja magnetyczna''.

Obrócenie ruchomego elementu o kąt $ \alpha $ spowoduje że będzie on próbował powrócić do położenia $ \alpha = 0 $, w którym reluktancja obwodu jest najmniejsza. W silnikach VRM przepływ prądu stałego przez uzwojenia powoduje, że zęby wirnika ustawiają się naprzeciw zasilanego uzwojenia.

\insertImg{grafiki/vrm_przekroj.png}
		{Silnik o zmiennej reluktancji (VRM)- Przekrój.}
		{przepiorkowski}

W odróżnieniu od VRM, wirnik silnika krokowego PM ({\em Permanent Magnet}- Magnes trwały) nie posiada zębów. Jest wykonany w postaci walca naprzemiennie namagnesowanego biegunami N i S. Specyficzny rodzaj zębów ma rdzeń stojana. \newline
W zależności od kierunku przepływu prądu w uzwojeniu przyciągane są odpowiednie bieguny wirnika (resunek \ref{grafiki/pm_stepper.png}).

\insertImg{grafiki/pm_stepper.png}
		{Zasada działania silnika PM.}
		{ni}
		
Główną zaletą silnika krokowego PM jest zastosowanie magnesów trwałych w stojanie, dzięki czemu nie ma potrzeby stosowanie szczotek (jak w silnikach DC) oraz jego niska cena. Wadą tego typu maszyny jest relatywnie niski moment obrotowy i brak możliwości rozwijania dużych prędkości obrotowych. \\

Silnik hybrydowy łączy w sobie cechy obu rozwiązań, dzięki czemu zostały poprawione takie parametry jak:
\begin{easylist}
	& moment obrotowy, 
	& maksymalna prędkość obrotowa, 
	& rozdzielczość kroku.
	\\
\end{easylist}

Niestety silniki hybrydowe są około 2-3 razy droższe od silników PM. Wirnik silnika HB ({\em Hybrid Motor}) jest zbudowany z uzębionych nabiegunników i magnesu trwałego, powodującego naprzemienne magnesowanie zębów biegunami N i S. Uzębiony stojan konstrukcją przypomina ten z silnika VRM. 

\insertImgSetSize{grafiki/stepper_budowa.png}
		{70}
		{Konstrukcja dwufazowego silnika krokowego hybrydowego.}
		{enggarage}
		
W silniku HB wirujące pole stojana obraca wirnik z jednego położenia do drugiego na zasadzie jak w silniku VRM. Jest to możliwe dzięki przesunięciu ,,północnej'' i ,,południowej'' części wirnika o pół ząbka. Silniki HB dzięki poprawionym parametrom są obecnie najbardziej popularnymi silnikami krokowymi, pomimo niskiej ceny silników PM.
		
\insertImgSetSize{grafiki/stepper_rotor.png}
		{70}
		{Budowa wirnika silnika krokowego hybrydowego dwufazowego.}
		{edw_sierpien_2002}

\subsubsection{Sterowanie}
\label{sss:sterowanie_krokowy}

W pierwszej kolejności należy rozróżnić sposób zasilania silnika krokowego, który może być unipolarny lub bipolarny (rys. \ref{grafiki/stepper_zasilanie_uni_bi.png}).

\insertImg{grafiki/stepper_zasilanie_uni_bi.png}
		{Sposób sterowania (zasilania) silnika krokowego dwufazowego- z lewej unipolarnie, z prawej bipolarnie.}
		{hackaday}

Przy zasilaniu unipolarnym, dla pojedynczego uzwojenia na zewnątrz silnika wyprowadzone są trzy przewody- Oba końce cewki i odczep w połowie jej długości. Podłączając do odczepu ,,+'' zasilania można sterować połówkami uzwojeń zwierając w odpowiedniej kolejności i na odpowiedni czas końcówki do masy. W zasilaniu bipolarnym, prąd płynie zawsze przez całe uzwojenie (z silnika wyprowadzone są tylko dwa przewody na cewkę). Zaletą wariantu unipolarnego jest możliwość znacznego uproszczenia układu sterowania. Poważną wadą natomiast jest to, że prąd podawany jest zawsze tylko na połowę danego uzwojenia, co ma negatywne przełożenie na moment obrotowy. 

\insertImg{grafiki/stepper_simple.png}
		{Uproszczony model silnika krokowego dwufazowego o kroku 90\degree.}
		{howtomecha}

Oddzielną sprawą jest sposób/ sekwencja podawania impulsów prądowych do uzwojeń silnika. Sterowanie pod tym względem dzieli się na:

\begin{easylist}
	& Falowe, 
	& Pełnokrokowe, 
	& Półkrokowe,
	& Mikrokrokowe.
	\\
\end{easylist}

Do pomocy w wyjaśnieniu poszczególnych typów sterowania posłuży rysunek \ref{grafiki/stepper_simple.png}, prezentujący uproszczony silnik krokowy. Maszyna składa się z wirnika zbudowanego jak na rys. \ref{grafiki/stepper_rotor.png} tyle że z ilością zębów (na wirniku i stojanie) pozwalającą osiągnąć pełen obrót w czterech pełnych komutacjach. Silnik posiada 2 uzwojenia A i B nawinięte w taki sposób, że pary A/A' i B/B' są ze sobą połączone (uzwojenie podzielone na dwie sekcje). Przykładając napięcie do cewki A, w zależności od polaryzacji tego napięcia, prąd będzie wpływał przez uzwojenie A i wypływał z uzwojenia A' lub na odwrót. Tak samo ma się sytuacja w przypadku cewki B.

\tocLessLeftNorm{Sterowanie falowe}

W przypadku braku zasilania, wirnik silnika ustawi się w jednej z czterech spoczynkowych pozycji (0\degree, 90\degree, 180\degree lub 270\degree). Dzieje się tak dlatego, że namagnesowany wirnik usiłuje zająć taką pozycję, aby reluktancja obwodu magnetycznego była jak najmniejsza, tzn. aby namagnesowane zęby wirnika były ustawione w jednej linii z dowolną parą zębów stojana. Znikomy moment obrotowy powodujący takie pozycjonowanie to inaczej moment bezprądowy (spoczynkowy). 

\insertImgSetSize{grafiki/stepper_sterowanie_falowe.eps}
		{100}
		{Sekwencja zasilania uzwojeń przy sterowaniu falowym (obrót wirnika z rys. \ref{grafiki/stepper_simple.png} według wskazówek zegara).}
		{oprWlasne}

Rysunek \ref{grafiki/stepper_sterowanie_falowe.eps} prezentuje sekwencję dzięki której wirnik obróci się o trzy pełne obroty w prawo przy zastosowaniu sterowania falowego. Na osi odciętych znajdują się kolejne chwile w czasie, a na osi rzędnych wartość potencjału przyłożonego do danego zacisku (wartości 1 i 0 prezentujące stan wysoki lub niski potencjału). Aby wirnik obracał się w przeciwnym kierunku, sekwencję należy odtworzyć w odwrotnej kolejności. Jak widać na przebiegach, w danym momencie w czasie prąd znajduje się zawsze tylko w jednym z uzwojeń, co nie umożliwia uzyskania maksymalnego dla danej maszyny momentu obrotowego. Pomimo że sterowanie falowe nie zapewnia uzyskania optymalnego momentu obrotowego, to przy zastosowaniu w niektórych aplikacjach jest wystarczające, a ponadto jest jest najprostsze do zaimplementowania.

\clearpage

\tocLessLeftNorm{Sterowanie pełnokrokowe}

W trybie sterowania pełnokrokowego prąd płynie przez oba uzwojenia jednocześnie. Nie ma to wpływu na prędkość obrotową wirnika, ale ma znaczący wpływ na moment obrotowy, który jest dwa razy większy niż w trybie falowym. Pełny obrót wirnika dalej jest osiągany w czterech "ruchach", ale absolutna pozycja kątowa w każdym z nich jest względem trybu falowego przesunięta o 45\degree.

\insertImgSetSize{grafiki/stepper_sterowanie_krokowe.eps}
		{100}
		{Sekwencja zasilania uzwojeń przy sterowaniu pełnokrokowym (obrót wirnika z rys. \ref{grafiki/stepper_simple.png} według wskazówek zegara).}
		{oprWlasne}
		
Rysunek \ref{grafiki/stepper_sterowanie_krokowe.eps} przedstawia przebiegi dla 3 obrotów, przy sterowaniu pełnokrokowym. Kolejne pozycje wirnika są następujące: 45\degree w pierwszym interwale czasowym, 135\degree w drugim, 225\degree w kolejnym i 315\degree w ostatnim (dla danego obrotu). Potem cykl się powtarza. Tryb pełnokrokowy pozwala osiągnąć optymalny dla danego silnika moment obrotowy, ale niestety rozdzielczość obrotu jest ograniczona do ilości kroków na obrót danego modelu maszyny. Typowe wartości kroków na obrót w silnikach krokowych hybrydowych to 100 (krok o rozdzielczości 3.6\degree), 200 (1.8\degree) i 400 (0.9\degree). Aby uzyskać większą rozdzielczość kroku, należy zastosować sterowanie półkrokowe lub mikrokrokowe.

\clearpage

\tocLessLeftNorm{Sterowanie półkrokowe}

Tryb półkrokowy jest połączeniem sterowania falowego i pełnokrokowego. Rozdzielczość obrotowa względem poprzednich trybów jest dwa razy większa, dlatego że uzwojenia naprzemiennie zasilane są pojedynczo/ oba jednocześnie.

\insertImgSetSize{grafiki/stepper_sterowanie_polkrokowe.eps}
		{100}
		{Sekwencja zasilania uzwojeń przy sterowaniu półkrokowym (obrót wirnika z rys. \ref{grafiki/stepper_simple.png} według wskazówek zegara).}
		{oprWlasne}
		
Wadą takiego rozwiązania jest to że w co drugim kroku moment obrotowy będzie mniejszy. Może to spowodować szarpanie wirnika przy dużych prędkościach obrotowych. Przy niskich prędkościach i obciążeniach natomiast, efekt ten rzadko występuje i można uzyskać polepszoną płynność ruchu, w porównaniu do poprzednich trybów. Aby uniknąć efektu zmiennego momentu, można zastosować sterowanie półkrokowe ze zmiennym prądem (rys. \ref{grafiki/stepper_sterowanie_polkrokowe_zmienne.eps}). Na rysunku widać że w momentach kiedy różnica potencjałów występuje tylko na jednym z uzwojeń, jest ona dwa razy większa niż w przypadku kiedy zasilane są obie cewki. Prąd płynący w uzwojeniach jest proporcjonalny do przyłożonego do nich napięcia, dlatego moment się wyrównuje i silnik nie jest narażony na szarpanie. Rozwiązanie te ma jednak wady- Zakładając że dla momentu w czasie, w którym przez oba uzwojenia płynie prąd, wartość tego prądu jest maksymalna (nominalna) dla danego modelu silnika w każdym z uzwojeń, oznacza to że w chwilach kiedy kiedy napięcie jest podawane na tylko jedno uzwojenie, prąd w nim płynący jest większy od nominalnego (tak jak na rys. \ref{grafiki/stepper_sterowanie_polkrokowe_zmienne.eps}). To naraża silnik na nadmierne grzanie, a nawet na uszkodzenie. Z drugiej strony można założyć że nominalny prąd ma płynąć w momencie zasilania jednego uzwojenia, a mniejszy przy zasilaniu dwóch. Ten tryb jest bezpieczny, ale optymalny moment obrotowy nie jest uzyskiwany.

\insertImgSetSize{grafiki/stepper_sterowanie_polkrokowe_zmienne.eps}
		{100}
		{Sekwencja zasilania uzwojeń przy sterowaniu półkrokowym ze zmiennym prądem (obrót wirnika z rys. \ref{grafiki/stepper_simple.png} według wskazówek zegara).}
		{oprWlasne}
		
\tocLessLeftNorm{Sterowanie mikrokrokowe}

Największą rozdzielczość ruchu (teoretycznie nieograniczoną, a praktycznie ograniczoną możliwościami układu sterującego) można uzyskać przy zastosowaniu trybu mikrokrokowego. Aby obrócić wirnik z rys. \ref{grafiki/stepper_simple.png} w najpłynniejszy możliwy sposób, kształt prądu podawany na uzwojenia powinien być naturalny dla danej konstrukcji silnika, czy też dla ruchu który wykonuje element ruchomy- W tym wypadku jest to ruch obrotowy. 

\begin{equation} \label{eq:stepper2}
	\sin^2\theta + \cos^2\theta = 1
\end{equation}

Równanie \ref{eq:stepper2} prezentuje wzór na jedynkę trygonometryczną. Można go użyć, aby uzyskać dowolną pozycję wirnika pomiędzy kolejnymi naturalnymi ułożeniami wirnika. W tym momencie należy rozróżnić dwa rodzaje pozycji kątowych wirnika silnika krokowego- magnetyczną oraz fizyczną. W przypadku silnika z rys. \ref{grafiki/stepper_simple.png} obie pozycje będą się zawsze pokrywać, ze względu na to że potrzeba jedynie czterech pełnych kroków aby uzyskać pełen obrót- kąt ,,magnetyczny'' 45\degree odpowiada obrotowi wirnika o 45\degree. Dla rzeczywistej sytuacji, w której silnik posiada np. 200 kroków na obrót, jeden stopień magnetyczny odpowiadałby już jedynie $ \frac{1}{50} $ stopnia w obrocie. Innymi słowy potrzeba 50 obrotów magnetycznych na 1 obrót fizyczny wirnika.

\insertImg{grafiki/circle_sine.eps}
		{Ćwiartki okręgu i funkcji $ \sin x $ w kartezjańskim układzie współrzędnych}
		{wikimedia}
		
Tak więc, podając na uzwojenia silnika prąd o kształcie funkcji {\it sinus} (na pierwsze uzwojenie) i {\it cosinus} (na drugie), suma ich kwadratów powinna zawsze wynosić jeden. 

\insertImgSetSize{grafiki/stepper_sin_cos.eps}
		{85}
		{Kształt prądów podawanych na uzwojenia A i B w celu uzyskania optymalnej płynności obrotu w silniku krokowym hybrydowym dwufazowym.}
		{oprWlasne}
		
Prezentowane przebiegi z rys. \ref{grafiki/stepper_sin_cos.eps} obrócą wał silnika o 360\degree magnetycznych. Maksymalna amplituda po znormalizowaniu reprezentuje wartość natężenia nominalnego prądu dla danego modelu silnika. Zazwyczaj silniki krokowe są sterowane z przekształtników, które same są zasilane napięciem stałym. Rozdzielczości uzyskiwane w typowych sterownikach wynoszą od 8 do 32 mikrokroków na pełen krok.  \\

\insertImgSetSize{grafiki/stepper_sin_cos_8.png}
		{70}
		{Kształt prądów podawanych na uzwojenia A i B przy sterowaniu mikrokrokowym z rozdzielczością $ \frac{1}{8} $ kroku.}
		{lab_boni_holub}
		
Tabela \ref{tab:stepper_sterowanie} prezentuje podsumowanie omówionych trybów sterowania silnikami krokowymi dla zasilania bipolarnego i unipolarnego.

\insertTab{|c|c|c|c|c|c|c|c|c|}
{%
\hline
\makecell{Rodzaj\\sterowania} & \multicolumn{2}{|c|}{Falowe} & \multicolumn{2}{|c|}{Pełnokrokowe} & \multicolumn{2}{|c|}{Półkrokowe} & \multicolumn{2}{|c|}{Mikrokrokowe} \\ 
\cline{2-9} & Bi & Uni & Bi & Uni & Bi & Uni & Bi & Uni \\ \hline
\makecell{Moment\\obrotowy} & 50\% & 25\% & 100\% & 50\% & \makecell{50-\\100}\% & \makecell{25-\\50}\% & 50\% & 25\% \\ \hline
\makecell{Komplikacja\\sterownika} & średnia & mała & średnia & mała & duża & duża & duża & duża \\ \hline
\makecell{Płynność\\ruchu} & mała & mała & mała & mała & duża & duża & duża & duża \\ \hline
\makecell{Rozdzielczość\\kroku} & mała & mała & mała & mała & średnia & średnia & duża & duża \\ \hline
}
{Cechy poszczególnych trybów sterowania silnikiem krokowym}
{przepiorkowski}
{tab:stepper_sterowanie}

\subsubsection{Zalety i wady silników krokowych}

Możliwość sterowania silnikiem krokowym za pomocą impulsów jest dużą zaletą- Przy odpowiednio dobranym układzie napędowym, do pracy w ograniczonym obszarze dokładności silnik krokowy nie potrzebuje żadnego sprzężenia zwrotnego (praca w otwartej pętli sprzężenia zwrotnego). Teoretycznie pozycja wirnika jest znana po każdej komutacji wykonanej przez układ sterujący. Niestety, wirnik i jego obciążenie posiadają pewną bezwładność co sprawia że przed osiągnięciem ustalonej pozycji występują drgania. \\
		
W przypadkach kiedy wymagana jest duża precyzja sterowania, tak duże drgania są bardzo niekorzystne. Wybrane sposoby eliminacji tych problemów to:

\begin{easylist}
	& Tłumienie mechaniczne poprzez zwiększenie momentu tarcia.
	& Sprzęgła elastyczne i przekładnie paskowe zastosowane w układzie napędowym, w celu redukcji przenoszonych drgań do obciążenia.
	& Wybór pracy półkrokowej lub mikrokrokowej- Im bardziej przebieg prądu jest zbliżony do sinusoidalnego, tym ciszej i stabilniej pracuje silnik krokowy.
	& Zastosowanie silnika krokowego o mniejszym skoku.
	\\
\end{easylist} 

\insertImgSetSize{grafiki/stepper_drgania.png}
		{60}
		{Wykres czasowy ilusytujący położenie wirnika podczas wykonywania skoku}
		{edw_sierpien_2002}

Należy pamiętać o wielkości i częstotliwości oscylacji decyduje głównie masa wirnika oraz wielkość i rodzaj obciążenia.





\subsection{Silniki prądu stałego (DC)}

\subsubsection{Budowa}

\subsubsection{Zasada działania}

\subsubsection{Sterowanie}

\clearpage


\subsection{Silnik bezszczotkowy z wirującym magnesem (BLDC)}

Konstrukcja silnika BLDC ({\it Brushless DC Motor}) jest odwrotnością silnika komutatorowego z magnesem trwałym: uzwojenia znajdują się w stojanie, a wirnik jest wykonany z odpowiednio ukształtowanego magnesu. Maszyny te są zaliczane do grupy silników prądu stałego, ale w praktyce pradem stałym zasilany jest jedynie falownik sterujący silnikiem. 

\insertImgSetSize{grafiki/bldc_przekroj.jpg}
		{50}
		{Przekrój typowego silnika PM BLDC: a) z wirnikiem wewnętrznym, b) z wirnikiem zewnętrznym, 1- jarzmo stojana (twornika), 2- zęby stojana, 3- żłobek z uzwojeniami, 4- wirnik, 5- magnes trwały}
		{krykowski}
		
Silniki BLDC są konstrukcyjnie bardzo podobne do silników krokowych- można przyjąć że silniki BLDC są pewną podgrupą silników krokowych o kroku 180\degree (silniki 2-fazowe) lub 120\degree (3-fazowe). Kryterium odróżniające silniki krokowe od BLDC może być takie, że w większości aplikacji dla poprawnej silnika BLDC jego sterownik musi dysponować sygnałem zwrotnym o aktualnym położeniu kątowym wirnika (np. z czujników Halla). Silniki krokowe natomiast mogą pracować w otwartej pętli sprzężenia zwrotnego.

\insertImgSetSize{grafiki/bldc_gwiazda.jpg}
		{50}
		{Konfiguracja uzwojeń trójfazowego silnika BLDC w gwiazdę}
		{przepiorkowski}
		
Uzwojenia silnika BLDC dla maszyny 3-fazowej są układane w gwiazdę (rys. \ref{grafiki/bldc_gwiazda.jpg}). Prąd z układu sterującego jest doprowadzany do wyprowadzeń A1, B1 i C1 podczas gdy pozostałe końce uzwojeń (A2, B2, C2) są ze sobą zwarte. Oznacza to że w tym rodzaju maszyny, prąd przepływa zawsze przez co najmniej dwa uzwojenia, w jednym z dwóch możliwych kierunków (np. kiedy wyższy potencjał jest na C1, a niższy na A1 wtedy od C1 do A1). Rysunek \ref{grafiki/bldc_sterowanie.eps} przedstawia tą zależność.

\insertImgSetSize{grafiki/bldc_sterowanie.eps}
		{100}
		{Uproszczone przebiegi prądów w uzwojeniach silnika 3-fazowego sterowanego bipolarnie}
		{oprWlasne}
		
Zmiany momentu obrotowego w funkcji kąta obrotu są w silniku trójfazowym niewielkie i można je jeszcze bardziej zredukować stosując impulsy prądowe o kształcie trapezowym. 
Tak jak zaprezentowano na rys. \ref{grafiki/bldc_przekroj.jpg} silniki bezszczotkowe mogą mieć różną konstrukcję wirnika. W przypadku budowy klasycznej wirnik w kształcie walca jest umieszczony wewnątrz stojana z uzwojeniami. W silniku z wirnikiem zewnętrznym magnetyczny wirnik ma kształt kubka i obraca się wokół nieruchomego stojana. \\

Do głównych zalet silników BLDC należy zaliczyć to że zjawisko rezonansu mechanicznego jest o wiele mniejsze niż w przypadku silników krokowych, oraz mogą one uzyskiwać prędkości obrotowe o wiele większe te pierwsze (powyżej 10000 RPM). Ponadto, przy tej samej wartości momentu obrotowego silnik BLDC może być mniejszy i lżejszy od krokowego. Z racji tego że BLDC nie posiada komutatora, a co za tym idzie szczotek, jego żywotność jest bardzo długa i jest praktycznie ograniczona żywotnością mechaniki układu napędowego. Największą niedogodnością jest natomiast to, że z silnikiem bezszczotkowym z wirującym magnesem nie da się w praktyce pracować w otwartej pętli sprzężenia zwrotnego.

\clearpage



% --------------------------------------------------------------------
% Bibliografia
% ustawienie bibliografii
\nocite{*}
\bibliographystyle{acm-i}
\bibliography{bib_database} 
\clearpage

\listoftables
\clearpage

\listoffigures
\clearpage

\lstlistoflistings




\end{document}

